\chapter{Literature Review}

\section{Introduction to the Literature Review}

This chapter examines the current state of research on event-centric knowledge graphs and their evaluation, establishing the theoretical foundations and identifying gaps that motivate the present work. The review is organized thematically, beginning with conceptual foundations that define core terminology and theoretical frameworks, then examining prior work on EKG construction and evaluation, analyzing methodological patterns across the literature, and finally synthesizing identified gaps that directly inform this research's objectives.

The structure reflects the multidisciplinary nature of EKG evaluation, drawing from knowledge graph quality assessment, temporal information extraction, causal reasoning, and semantic web technologies. By critically comparing approaches across these areas, the review reveals a persistent fragmentation: different research communities evaluate different aspects of event-centric knowledge using incompatible metrics and benchmarks, leaving no comprehensive framework for holistic EKG quality assessment.

\section{Conceptual and Theoretical Foundations}

\subsection{Event-Centric Knowledge Graphs}

Event-centric knowledge graphs represent a fundamental shift from traditional entity-centric knowledge representation. While conventional knowledge graphs organize information around static entities and their relationships—such as ``Person A works at Company B''—EKGs place events as first-class nodes and structure surrounding information around those events \cite{sem2011}. This design choice reflects the observation that many real-world phenomena are inherently temporal and change-oriented, requiring representation of not just what entities exist, but what happened, when it happened, who was involved, and how events relate to each other.

The Simple Event Model (SEM) provides a foundational ontology for event representation, designed with minimal semantic commitment to maintain interoperability across domains while providing usable schema for actors, places, and temporal anchors \cite{sem2011}. SEM's flexibility is powerful but introduces evaluation challenges: ``completeness'' and ``correctness'' become partly context-dependent and fit-for-purpose rather than absolute measures. An event representation might be complete for timeline generation but insufficient for causal reasoning, depending on which contextual attributes and relationships are captured.

\subsection{Temporal Modeling and Constraints}

Temporal modeling adds significant complexity to event representation and evaluation. Events can be anchored explicitly through timestamps or only inferentially through temporal relations; temporal constraints can be underspecified; and logically equivalent temporal structures may be expressed through different edge sets. TimeML emerged as an influential standard precisely because it offers rich specification for events, temporal expressions, and ordering relations, highlighting that temporal evaluation cannot reduce to simple timestamp matching but must reason over structured temporal constraints \cite{timeml2003}.

This temporal complexity has direct implications for evaluation: two EKGs might represent the same temporal reality through different but logically consistent structures, yet traditional precision-recall metrics would penalize such differences. Temporal closure—the set of all relations logically implied by explicitly stated relations—becomes essential for fair evaluation, as demonstrated in temporal information extraction research \cite{tempeval2011}.

\subsection{Evaluation Dimensions for Knowledge Graphs}

The broader knowledge graph quality literature establishes that ``quality'' is inherently multidimensional. Surveys on KG quality control identify multiple evaluation dimensions including accuracy, completeness, consistency, timeliness, and provenance \cite{kgquality2021}. More recent comprehensive surveys emphasize that quality management must span assessment, detection, correction, and completion, particularly for dynamic graphs operating under open-world assumptions where missing information does not imply falsity \cite{kgquality2024}.

For EKGs specifically, quality assessment must address not only correctness of extracted triples but also whether the graph faithfully represents event structure (what happened), event grounding (who/where/when), and event logic (how events relate temporally and causally). This multifaceted nature of EKG quality motivates the need for evaluation frameworks that integrate multiple dimensions rather than focusing on isolated aspects.

\section{Prior Work on Event-Centric Knowledge Graphs}

\subsection{Construction of Event-Centric Knowledge Graphs}

The ecosystem of event-centric knowledge graphs reflects diverse construction approaches and application goals. EventKG represents a prominent integration-focused approach, combining event information from multiple reference sources including Wikipedia, Wikidata, and DBpedia to create a large-scale, multilingual event knowledge base \cite{eventkg2018}. EventKG emphasizes canonical representation, provenance tracking, and comprehensive temporal-spatial coverage, positioning its contribution as improving completeness compared to individual source datasets.

ASER (Activities, States, Events, and their Relations) takes a different approach, focusing on extracting eventuality knowledge from massive text corpora using syntactic pattern extraction and a typed relation inventory influenced by discourse relation schemes \cite{aser2019}. Rather than integrating structured sources, ASER emphasizes scale and coverage of commonsense event knowledge extracted from natural language.

The Open Event Knowledge Graph (OEKG) emphasizes cross-lingual access and multi-dataset integration, assembling seven datasets spanning question answering, entity recommendation, and named entity recognition through named-graph-based integration and linking to EventKG \cite{oekg2023}. OEKG's design prioritizes practical use cases such as hybrid question answering over knowledge graphs and news, and language-specific event recommendation.

Domain-specific EKGs address specialized requirements. NewsReader-style event-centric knowledge graphs focus on building temporal event networks from news corpora, explicitly modeling information related to events to provide temporal dimensions for narrative understanding \cite{newsreader2016}. Geographic knowledge graphs like AugGKG (grid-augmented geographic knowledge graph) emphasize spatio-temporal query support and efficiency for location-based event analysis \cite{auggkg2023}.

\subsection{Evaluation of Knowledge Graphs}

\subsubsection{Intrinsic Evaluation Approaches}

Intrinsic evaluation assesses knowledge graph quality through internal properties without reference to downstream tasks. Constraint and schema validation provides a foundational layer: W3C SHACL (Shapes Constraint Language) offers a standard for validating RDF graphs against shape constraints, formalizing syntactic and structural validity checks including cardinalities, datatypes, required properties, and link patterns \cite{shacl2017}. Test-driven linked data quality assessment extends this approach by framing quality checks as reusable test cases, often expressed as SPARQL queries, inspired by software testing methodologies \cite{testdriven2014}.

Frameworks like Luzzu institutionalize the ``metric plugin + scalable assessment'' approach for linked data quality, emphasizing extensibility and streaming evaluation to handle large-scale graphs \cite{luzzu2016}. These automated validation approaches provide repeatable baselines but cannot capture semantic correctness or fitness for specific reasoning tasks.

\subsubsection{Extrinsic Evaluation Approaches}

Extrinsic evaluation measures knowledge graph quality through downstream task performance. The KGrEaT framework explicitly advocates this approach, evaluating KGs via tasks such as classification, clustering, and recommendation using fixed experimental setups and entity mapping to compare graphs as inputs \cite{kgreat2023}. KGrEaT highlights that accessibility features like labels and descriptions can heavily influence downstream utility, suggesting that intrinsic completeness measures alone may not predict practical effectiveness.

For EKGs specifically, task-based evaluation often focuses on event-centric applications. Event-QA provides a benchmark for complex SPARQL queries over EventKG with multilingual verbalizations, enabling question-answering-based evaluation that probes whether event-entity relations and temporal constraints support query answering \cite{eventqa2020}. ChronoGrapher demonstrates utility through user studies on event-centric prompting for question answering, reporting improved answer groundedness when prompts are enriched with event-centric KG triples compared to generic entity-centric triples \cite{chronographer2023}.

\subsubsection{Scalable Human-in-the-Loop Auditing}

Given the infeasibility of complete manual annotation for large EKGs, sampling-based auditing with statistical guarantees has emerged as a practical approach. Research on efficient KG accuracy evaluation introduces sampling designs including stratified sampling and reservoir-style incremental approaches that dramatically reduce human annotation costs while producing meaningful accuracy estimates \cite{efficienteval2019}. These methods are particularly relevant for evolving EKGs where continuous quality monitoring is needed but full re-annotation is impractical \cite{incrementaleval2019}.

\subsection{Frameworks and Toolkits for Evaluation}

Existing evaluation frameworks provide partial solutions but lack comprehensive EKG-specific coverage. General KG quality frameworks emphasize multidimensional assessment but do not address event-specific features like temporal density, causal coherence, or event evolution patterns. Temporal information extraction frameworks like tieval standardize corpus access and provide domain-specific operations including temporal closure and temporal awareness metrics, explicitly targeting comparability across datasets with incompatible annotation schemes \cite{tieval2023}. However, tieval focuses on temporal relation extraction rather than holistic EKG evaluation.

Recent work explores LLM-assisted validation for knowledge graphs. KGValidator proposes LLM-based methods for validating KG completion outputs to reduce manual labeling requirements, while other frameworks combine LLM and human feedback loops for KG optimization \cite{kgvalidator2024}. These approaches are relevant for EKGs given that event triples often exist at the boundary between structured constraints and linguistic ambiguity, but they introduce new challenges around model bias and judgment reliability.

\section{Methodological Patterns in the Literature}

\subsection{Common Evaluation Designs}

Across the EKG literature, evaluation methods cluster into several recurring patterns, each valuable but often partial in scope.

\subsubsection{Triple-Level Precision, Recall, and F1}

Triple-level precision, recall, and F1 scores against human annotation remain common when evaluating relationship extraction or event threading. Early work on event evolution graphs evaluates extracted evolution relationships using precision and recall against manually annotated relations \cite{eventevolution2009}. This paradigm persists because it is straightforward and comparable to classic information extraction evaluation, but it suffers from inherent limitations: annotation scope constraints, and the fact that equivalent high-level event structures can map to multiple valid edge sets.

\subsubsection{Coverage and Attribute-Level Correctness}

Coverage and completeness metrics dominate evaluations for integrated encyclopedic EKGs. EventKG illustrates this pattern by reporting coverage and completeness over event representations and comparing presence of time and location attributes across sources \cite{eventkg2018}. Later work evaluates event identification, time fusion, and location fusion using user annotations over sampled events, reporting correctness, missingness, and derived precision for conflicting temporal information and locations. While this validates key EKG primitives (eventness, dates, places), it leaves higher-order properties like causality chain quality or multi-hop reasoning support largely unmeasured.

\subsubsection{Crowdsourced and Expert Judgment}

Human judgment is frequently employed where gold standards are unavailable. NewsReader-style ECKGs explicitly acknowledge lacking proper gold standards and therefore rely on human judgment over randomly sampled event triples \cite{newsreader2016}. ASER uses formal crowdsourcing protocols, sampling extracted eventualities and asking annotators to judge whether extractions capture sentence meanings, and comparing ASER's knowledge to ConceptNet while evaluating utility on commonsense reasoning benchmarks \cite{aser2019}.

\subsubsection{Temporal-Specific Evaluation}

Temporal evaluation has evolved beyond simple precision-recall-F1, motivated by the structured nature of temporal constraints. UzZaman and Allen propose temporal evaluation using temporal closure to reward logically implied relations, arguing that traditional methods miss equivalences and fragment evaluation across subtasks \cite{tempeval2011}. QA TempEval advances this by recentering evaluation on temporal question answering performance rather than edge-level scoring, arguing that QA better reflects whether extracted temporal information supports end-user tasks \cite{qatempeval2015}.

\subsubsection{Causality Evaluation}

Causality evaluation in EKG contexts typically borrows from causality identification tasks over text using specialized corpora. Causal-TimeBank provides a TimeML-inspired annotation framework and dataset capturing temporal and causal relations for evaluating extraction systems \cite{causaltimebank2014}. The Event StoryLine Corpus targets story-level temporal and causal relation detection across document collections \cite{storyline2017}. Multilingual causality benchmarks like MECI and domain-focused shared tasks like FinCausal 2020 demonstrate expanding scope toward multilingual and domain-specific settings \cite{meci2022}, but these efforts remain largely separated from holistic EKG evaluation frameworks that connect causal edge correctness to graph-level inference behavior.

\subsubsection{Graph-Structure-Level Evaluation}

Graph-structure-level evaluation appears in event graph research but is not widely adopted in EKG practice. Glavaš and Šnajder propose overall quality metrics based on tensor products between automatically and manually constructed graphs, moving toward holistic structural similarity rather than edge-by-edge scoring \cite{glavas2014}. The recent STAGE benchmark introduces Event-Structure Consistency metrics that prioritize structural coherence and faithful abstraction over exact matching, explicitly acknowledging that event summaries should be evaluated non-symmetrically through coverage and faithfulness dimensions \cite{stage2026}.

\subsection{Methodological Weaknesses}

Several methodological weaknesses recur across the literature, motivating the design choices in this research.

\subsubsection{Fragmented Evaluation Across Incompatible Metrics}

Different EKG projects validate different aspects of event-centric knowledge using incompatible benchmarks and metrics. The temporal information extraction community's experience is instructive: tieval explicitly attributes limited comparability to heterogeneous annotation schemes, inconsistent formats, and divergent metrics across datasets \cite{tieval2023}. These conditions are magnified in EKGs because they combine multiple subtasks—event detection, argument extraction, temporal grounding, causal inference, and cross-document linking—each with its own evaluation traditions.

\subsubsection{Scarcity of Gold Standards at EKG Scope}

A recurring bottleneck is the scarcity of gold standards at full EKG scope. NewsReader-style ECKG evaluation highlights that when proper gold standards are unavailable, evaluation falls back to human judgment over small samples, providing useful spot checks but not fully characterizing graph-level reasoning behavior under complex multi-hop queries \cite{newsreader2016}. Even in large integrated resources like EventKG, evaluation focuses on specific pipeline steps and attribute-level precision, leaving causal correctness, event evolution structure, and inference robustness largely outside the measured space \cite{eventkg2018}.

\subsubsection{Causality as an Underaddressed Dimension}

Causality represents an especially important and difficult gap. Causal relations are often implicit, entangled with temporal order, and domain-dependent. Dedicated corpora help measure causal edge extraction but do not validate whether an EKG's causal subgraph supports reliable downstream inference such as counterfactual queries, intervention-like reasoning, or causal chain explanations \cite{causaltimebank2014}. Temporal evaluation work shows that logically equivalent structures can lead to misleading scores if evaluation does not account for closure and global consistency—an argument that applies even more strongly to causal chains where multi-hop implications matter \cite{tempeval2011}.

\subsubsection{Event Granularity and Contextual Richness}

Event granularity and contextual richness are often acknowledged but rarely standardized in evaluation. SEM's design motivation includes representing ``who/what/when/where'' with flexibility and minimal commitment, but that flexibility means ``completeness'' depends on intended queries and domain expectations \cite{sem2011}. EventKG reports temporal and spatial coverage, but contextual completeness would also encompass participant roles, event types, provenance, and viewpoint-dependent descriptions, which are not uniformly measured \cite{eventkg2018}.

\subsubsection{Multilingual and Cross-Domain Robustness}

Multilingual and cross-domain robustness remains difficult to evaluate holistically. OEKG assembles multilingual resources and demonstrates cross-lingual use cases, Event-QA provides multilingual verbalizations for event-centric QA, and MECI targets multilingual causality \cite{oekg2023,eventqa2020,meci2022}. Yet these are typically evaluated in their own silos—QA accuracy, use-case demonstration, or causality F1—rather than under a unified EKG quality model that jointly measures cross-lingual alignment, temporal-causal consistency, and query utility.

\subsubsection{Static Evaluation of Evolving Knowledge}

The evolving nature of event knowledge is underrepresented in evaluation protocols. Surveys of KG evolution distinguish static, dynamic, temporal, and event KGs and emphasize that timeliness and continuous updates are core to modern KGs \cite{kgevolution2023}. Yet many evaluation studies remain one-off snapshots rather than continuous monitoring with incremental re-evaluation. Sampling-based incremental evaluation methods for evolving KGs exist in the general KG literature but are not commonly integrated into EKG evaluation toolchains \cite{incrementaleval2019}.

\section{Identified Gaps and Research Opportunities}

The literature review reveals five critical gaps that directly motivate this research and inform its objectives.

\subsection{Gap 1: Absence of Comprehensive, Multidimensional Frameworks}

Current EKG evaluation is fundamentally fragmented. Different projects assess isolated features—extraction precision, temporal coverage, QA utility, causal relation accuracy—using incompatible metrics and benchmarks. No existing framework simultaneously evaluates the temporal, causal, contextual, structural, and semantic dimensions that collectively define EKG quality. This fragmentation prevents systematic comparison of EKG systems, obscures trade-offs between quality dimensions, and provides no guidance for prioritizing improvements.

\textbf{Research Opportunity:} Develop a conceptual framework that identifies and organizes critical dimensions of EKG quality, establishing relationships between evaluation criteria and defining how they contribute to overall effectiveness. This addresses RQ1 (What are the essential dimensions of quality for event-centric knowledge graphs, and how can they be systematically measured?).

\subsection{Gap 2: Lack of Standardized Metrics for Event-Specific Features}

While precision, recall, and temporal consistency have established definitions, critical event-specific features lack standardized metrics. Causal coherence, event granularity, contextual completeness, and reasoning capability are acknowledged as important but are not systematically measured. Existing metrics borrowed from traditional KG evaluation fail to capture event-specific phenomena like temporal density, event evolution patterns, or causal chain consistency.

\textbf{Research Opportunity:} Develop a comprehensive metric suite that integrates established measures with novel metrics for previously unassessed dimensions, providing formal definitions and computational procedures. This addresses RQ1 and enables the comparative analysis in RQ3 (How do different EKG construction approaches perform across multiple quality dimensions, and what trade-offs exist?).

\subsection{Gap 3: Insufficient Tools for Practical, Automated Assessment}

Existing evaluation approaches often require manual annotation, custom scripts, or specialized expertise, limiting their practical applicability. While frameworks like Luzzu provide extensible architectures for linked data quality and tieval standardizes temporal evaluation, no comprehensive toolkit exists for automated, multidimensional EKG assessment that practitioners can readily apply to their own knowledge graphs.

\textbf{Research Opportunity:} Implement an open-source evaluation toolkit that operationalizes the framework through SPARQL-based analysis, graph-theoretic measures, and semantic validation, providing automated assessment with detailed reporting. This supports all three research questions by enabling empirical validation and practical adoption.

\subsection{Gap 4: Limited Validation of Comprehensive Approaches}

Most evaluation studies focus on validating specific extraction or integration methods rather than validating evaluation frameworks themselves. It remains unclear whether comprehensive, multidimensional evaluation reveals quality differences that narrow approaches miss, or whether the added complexity provides actionable insights beyond traditional metrics.

\textbf{Research Opportunity:} Validate the framework through empirical evaluation of multiple EKG systems and synthetic datasets, demonstrating its ability to reveal quality differences across dimensions and comparing findings with traditional evaluation approaches. This directly addresses RQ2 (Can a comprehensive evaluation framework reveal quality differences across EKG systems that are missed by traditional, narrow evaluation approaches?).

\subsection{Gap 5: Weak Connection Between Intrinsic Quality and Task Utility}

The relationship between intrinsic quality measures (completeness, consistency, coverage) and extrinsic task performance (QA accuracy, timeline generation, causal reasoning) remains poorly understood for EKGs. KGrEaT demonstrates that accessibility features can heavily influence downstream utility for general KGs, but similar analysis for event-specific quality dimensions and event-centric tasks is lacking.

\textbf{Research Opportunity:} Analyze how different quality dimensions correlate with task performance and identify which dimensions are most critical for specific applications. While full extrinsic validation is beyond the scope of this thesis, the framework design should enable future research connecting intrinsic metrics to task utility, and initial validation should examine whether comprehensive assessment provides insights that guide system improvement.

\section{Chapter Summary}

This literature review has examined the current state of event-centric knowledge graph evaluation, revealing a field characterized by significant advances in EKG construction but fragmented and incomplete evaluation practices. The review established conceptual foundations including event-centric representation principles, temporal modeling complexities, and multidimensional quality frameworks. It surveyed major EKG systems (EventKG, ASER, OEKG) and their construction approaches, analyzed evaluation methods spanning intrinsic validation, extrinsic task-based assessment, and human-in-the-loop auditing, and examined existing frameworks and toolkits.

The analysis of methodological patterns revealed recurring evaluation designs—triple-level precision-recall, coverage metrics, crowdsourced judgment, temporal-specific evaluation, causality assessment, and graph-structure evaluation—each providing valuable but partial insights. Critical methodological weaknesses include fragmented evaluation across incompatible metrics, scarcity of gold standards at EKG scope, underaddressed causality, unstandardized treatment of event granularity and contextual richness, difficulties in multilingual and cross-domain evaluation, and static evaluation of inherently evolving knowledge.

Five critical gaps emerge from this analysis: (1) absence of comprehensive, multidimensional frameworks; (2) lack of standardized metrics for event-specific features; (3) insufficient tools for practical, automated assessment; (4) limited validation of comprehensive approaches; and (5) weak connection between intrinsic quality and task utility. These gaps directly motivate the research objectives and questions established in Chapter 1.

Given these gaps, the next chapter describes the research design and methodology used to develop and validate a comprehensive evaluation framework for event-centric knowledge graphs. The methodology addresses the identified weaknesses by integrating multiple evaluation dimensions, formalizing event-specific metrics, implementing practical assessment tools, and conducting empirical validation across multiple EKG systems and synthetic datasets.
